\begin{frame}[t]{Freiburg Ion Trapping Group}
    \picpos{\linewidth}{0cm}{0cm}{group_slide.pdf}
\end{frame}

\begin{frame}[t]{Atom-Ion interaction}
    The ion induces a dipole moment in the atom ($p \propto \alpha E(r) \propto \frac{1}{r^2}$) that leads to
    charge-dipole interaction ($\propto\frac{p}{r^2}$)
   \[ 
        \setlength{\abovedisplayskip}{2pt}
        \setlength{\belowdisplayskip}{2pt}
       V_\text{atom-ion}(r) = \frac{C_6}{r^6} - \frac{C_4}{r^4} + \frac{\hbar \,l\,(l+1)}{\mu r^2}
       \]
    \picpos{0.45\linewidth}{.5\linewidth}{0cm}{BaLi_PECs_B.pdf}
    \picpos{0.3\linewidth}{1cm}{0cm}{atom_ion_potential.png}
\end{frame}

\begin{frame}[t]{Feshbach Resonances}
    Scattering state (open channel) can couple to a bound state (closed channel) leading to a pole in the
    scattering length $a$.
    \picpos{.4\linewidth}{.5cm}{0.5cm}{feshbach_resonance.pdf}
    \only<2>{
        \picpos{.46\linewidth}{.51\linewidth}{2.4cm}{feshbach_resonance_isolated.pdf}
        \begin{textblock*}{.5\linewidth}(.51\linewidth, 0.2cm)
           \small
           Magnetic fields can be used to shift potentials relative to each other (magnetic FR)
                   \[ a = a_\text{bg} \left( 1 - \frac{\Delta}{B-B_0} \right) \]
        \end{textblock*}
   }
\end{frame}
